\documentclass{article}
\title{身のまわり雑学帖 — なぜ?に答える読み物}
\author{TeX 図書館}

\begin{document}

\section{読み方 — 雑学は入口、理解は本編}

この本は教科書ではありません。\textbf{身のまわりの「なぜ?」を、ちゃんと答える読み物}です。ただし、1 つの「なぜ」を 1 行で終わらせません。\textbf{短い答え → しくみの説明 → 応用と深掘り → 図書館の該当箇所への扉}の 4 段構成で、雑学として読んでも、学習の入口としても使えるようにします。

各テーマの最後にある「続きはこちら」の案内先が、この図書館の本編です。雑学で興味を持った話は、本編で体系的にやり直すと記憶に残ります。統計の誤用事例集と同じく、読み物として全編通しで読んでも、好きな節だけ拾い読みしても構いません。

\section{空のなぜ}

\subsection{空はなぜ青い}

\textbf{短い答え}: 光は波で、青い光ほど空気の粒で強く散らされるから。

\textbf{しくみ}: 太陽光は 7 色の光の集まりで、波長が違います(赤が長波長、青が短波長)。空気の分子の大きさは光の波長よりずっと小さく、この小さな粒に当たった光は、\textbf{波長の 4 乗に反比例する強さ}で散乱します(レイリー散乱)。4 乗というのは急激な割合で、青の散乱は赤の約 5 倍。空を見上げると、太陽の光から散らされた青い光が目に入るので、空全体が青く見えます。

\textbf{応用と深掘り}: 夕焼けが赤いのは、同じ仕組みの裏面です。夕方は光が大気を斜めに長く通るので、青は途中で散らし尽くされて、散り残った赤が目に届く。\textbf{月には大気がないので、空は昼でも黒い}。散乱がなければ空の色はない、というのが正しい言い方です。

続きはこちら: 光の波としての性質は \textbf{高校物理やり直し}(音と光)で、散乱と色の話は \textbf{音楽理論入門}(音の物理)の対応物ごとに読めます。

\subsection{雪の結晶はなぜ六角形}

\textbf{短い答え}: 水の分子が結晶を作るとき、六角形の並びが最も安定だから。

\textbf{しくみ}: 水分子(H\(_{2}\)O)は、酸素の周りに水素が 105 度くらいの角度で付いた折れ曲がった形です。氷の結晶では、この形が決まった角度で繰り返され、\textbf{六角形のネットワーク}になります。結晶は成長しても分子の並びのパターンを保つので、どんなに枝分かれしても基本の 6 回対称を守ります。枝の形がぞれぞれ違うのは、成長中の温度と湿度の微細な違いに敏感に反応するためです。

\textbf{応用と深掘り}: 結晶の形が分子の並びで決まる、という見方は化学全体の入口です。塩が立方体、石英が六角柱。結晶構造の話は \textbf{高校化学やり直し}(粒子モデル)から始まり、\textbf{数値計算入門}(シミュレーション)で成長の過程を計算機で追えます。

\subsection{虹はなぜ弧}

\textbf{短い答え}: 雨粒の中で光が屈折と反射をし、太陽の反対方向の 42 度に強く出るから。

\textbf{しくみ}: 雨粒に入った光は、粒の中で 1 回反射して出ていきます。このとき入射と出射の角度は同じですが、\textbf{波長ごとに屈折の角度がわずかに違う}(分散)ので、色が分かれます。計算すると、光が最も強く出る方向は太陽の反対方向から 42 度。地面に立つ私から見ると、太陽の反対方向を中心に 42 度の円錐の縁が、色分けされた弧として見えます。

\textbf{応用と深掘り}: 光が媒質で曲がる(屈折)は、レンズ・眼鏡・カメラ・光ファイバーすべての原理です。全反射で光を端から端まで運ぶ光ファイバーは、\textbf{高校物理やり直し}(音と光)の応用編です。

\section{水と温度のなぜ}

\subsection{塩を撒くと雪はなぜ溶ける}

\textbf{短い答え}: 塩が溶けた水は、0 度でも凍らなくなる(凝固点が下がる)から。

\textbf{しくみ}: 純水は 0 度で凍ります。しかし塩を溶かすと、水の分子が氷の結晶を作るのを塩のイオンが邪魔します。\textbf{凍るには「純水 0 度」より低い温度まで冷えないと、結晶が完成しない}。これを凝固点降下と呼び、溶けた物の粒子数だけで決まります(種類は関係しません)。融雪剤は、この性質で氷点下でも水を保つ仕組みです。

\textbf{応用と深掘り}: 同じ原理が、冬の道路の融雪、アイスクリーム作り(塩と氷で 0 度以下を作る)、自動車の不凍液です。溶液の性質は \textbf{高校化学やり直し}(溶液の濃度・酸と塩基)で体系的に学べます。

\subsection{鍋はなぜ吹きこぼれる}

\textbf{短い答え}: 沸騰の前に鍋底が局所的に過熱され、泡が一気に成長するから。

\textbf{しくみ}: 理想的な沸騰では、泡が滑らかに発生します。しかし滑らかな鍋底には泡の核になる傷が少なく、底の水が沸点を超えても泡を作れずに過熱されます(過熱液)。ようやく泡ができたとき、その泡は\textbf{過熱された水を瞬時に蒸発させて急成長}し、鍋の中の水を押し上げる。これが吹きこぼれです。割り箸を入れると泡の核が増えて穏やかに沸騰するのは、核を人工的に増やすためです。

\textbf{応用と深掘り}: 過熱液の急激な沸騰は、電子レンジで長時間加熱したカップが手に取った瞬間に吹き出す事故の原因でもあります。状態変化の話は \textbf{高校化学やり直し}(粒子モデル)と \textbf{機械のしくみ入門}(冷蔵庫の節)の圧力と沸点の話につながります。

\subsection{冷凍焼けはなぜ起きる}

\textbf{短い答え}: 冷凍庫の中で、食材の水分が氷をすり抜けて昇華し、乾燥するから。

\textbf{しくみ}: 氷は 0 度以下でもゆっくり昇華します(固体から直接気体へ)。冷凍庫内は乾燥した環境なので、食材の表面の氷は数週間で昇華して失われ、白っぽく乾いた部分が残ります。これは腐敗ではなく\textbf{乾燥}で、味と食感は落ちますが安全性は保たれます。密閉袋に入れると、袋の中で飽和して昇華が止まるので防止できます。

\textbf{応用と深掘り}: 昇華は、凍結乾燥食品(インスタント食品の乾燥)の製造原理そのものです。氷を直接気化させて水を抜くと、形と風味を保ったまま乾燥できる。状態変化の 3 形態(融解・気化・昇華)は \textbf{高校化学やり直し}の粒子モデルの節で、温度と圧力の図とともに学べます。

\section{音のなぜ}

\subsection{遠くの雷はなぜゴロゴロ続く}

\textbf{短い答え}: 雷の音が山や建物・雲で何度も反射し、遅れて届く音が続くから。

\textbf{しくみ}: 雷の音は、空気中を約 340 m/s で伝わります(高校物理やり直し、音の節)。近くの雷は 1 つの経路の音が一瞬で届き「パンッ」と聞こえます。遠くの雷は、\textbf{雲・山・建物で反射した多数の経路}の音が、長さの違う遅れで次々に届くので、ゴロゴロと長く続いて聞こえます。また低い音ほど遠くまで届く(高い音は吸収されやすい)ので、遠雷ほど低く響きます。

\textbf{応用と深掘り}: 雷の距離は数えることができます。\textbf{光はほぼ瞬時、音は 340 m/s} なので、「光を見てから 3 秒ごとに約 1 km」。音の速さと反射の話は \textbf{高校物理やり直し}(音と光)の中心テーマで、超音波検査やソナーは同じ原理の応用です。

\subsection{救急車のサイレンはなぜ音が変わって聞こえる}

\textbf{短い答え}: 音源が近づくと波が圧縮されて高く、遠ざかると伸びて低く聞こえる(ドップラー効果)。

\textbf{しくみ}: 音は空気の密と疎の波です。音源がこちらに向かってくると、波の山の間隔(波長)が詰まって届くので、振動数が高く聞こえます。逆に遠ざかると間隔が広がり、低く聞こえる。\textbf{音源の振動数は変わっていないのに、観測される振動数が変わる}のがドップラー効果です。追い越す瞬間に「ピューピュー→ビャー」と音が下がるのは、接近から離反への切り替わりです。

\textbf{応用と深掘り}: 医療の超音波検査は、血球の流れがドップラーで音をずらすのを利用して血流を測ります。天体のスペクトルのずれ(宇宙の膨張の証拠)も同じ原理です。波の式 \( v = f\lambda \) とドップラーは \textbf{高校物理やり直し}(波の節)と \textbf{音楽理論入門}(音の物理)で詳しく扱います。

\subsection{無響室はなぜ怖い}

\textbf{短い答え}: 反射音がほぼゼロになり、普段は気づかない自分の体の音まで聞こえるから。

\textbf{しくみ}: 普段の部屋では、音が壁で反射して 0.1 秒単位で届き続けています(残響)。無響室は壁を楔形の吸音材で覆い、\textbf{反射音をほぼ消します}。すると耳は静寂ではなく、\textbf{自分の血流・鼓動・耳鳴り}のような体の中の音を拾い始めます。多くの人が数十分で落ち着かなくなるのは、脳が「音のない環境」に慣れていないからです。

\textbf{応用と深掘り}: 残響の制御は、コンサートホールの設計(反射音の長さで音の豊かさが変わる)の核心です。反射・共鳴・吸音の話は \textbf{高校物理やり直し}(音と光)と \textbf{音楽理論入門}(音の物理)に分けて収録しています。

\section{体のなぜ}

\subsection{指紋はなぜある}

\textbf{短い答え}: 皮膚の凸凹が、滑り止めと微細な感触のセンサーとして働くから。

\textbf{しくみ}: 指紋の凸凹は、物体に触れたときの\textbf{摩擦と振動}を増やします。凸部が先に触れて振動が生まれ、その振動のパターンが触覚受容器に、物の質感や滑りやすさの情報として伝わります。また汗が凸の先に分散して、湿り具合のバランスで摩擦を最適に保つ。胎児の頃に手の成長と皮膚の張力でできるため、\textbf{模様は生涯変わらず、人と完全には一致しない}。

\textbf{応用と深掘り}: 「模様が人ごとに違い、生涯変わらない」の 2 条件が、指紋認証の根拠です。スマートフォンの指紋センサーは、凸凹の起伏を電気的に読み取っています。感覚の仕組みは \textbf{心理学入門}(知覚の節)で、測る技術の一般論は \textbf{機械のしくみ入門}(センサーの節)につながります。

\subsection{高い音から聞こえなくなるのはなぜ}

\textbf{短い答え}: 内耳の高音を担当する細胞から、加齢と騒音で先に傷つくから。

\textbf{しくみ}: 耳の奥(蝸牛)には、振動数ごとに担当が分かれた感覚細胞が並んでいます。入口に近い側が高音を担当しますが、\textbf{高音の細胞は物理的に強い振動を受けやすい}位置にあり、加齢や大きな音で先に失われます。20 歳前後からゆっくり進み、聞こえる上限が年々下がる。若者だけが聞こえる高周波の着信音は、この加齢の差の利用です。

\textbf{応用と深掘り}: 聞こえ方の個人差は「内耳の何歳から壊れたか」の違いで、予防は騒音からの距離と時間の管理です。イヤホンの大音量が最も若いうちから効く、という点は実用上の最重要事項です。感覚が測られる仕組みは \textbf{心理学入門}(知覚)と \textbf{統計学入門}(測定の誤差)で学べます。

\subsection{めまいの正体は何}

\textbf{短い答え}: 内耳の三半規管が感じる回転の情報と、目や体の情報が食い違うから。

\textbf{しくみ}: 三半規管は、頭の回転を水の中の毛の揺れで検出するセンサーです。通常、目・体の感覚・三半規管の 3 つの情報は一致しています。車酔いや船酔いは、\textbf{体は静止を感じているのに、内耳は揺れを検出する}という不一致で、脳が混乱して吐き気が出ます。回り続けた後に立ちくらみがするのは、水が慣性で流れ続ける(センサーが遅れる)ためです。

\textbf{応用と深掘り}: 宇宙飛行士の宇宙酔いも同じ不一致です(重力がなく、体の感覚との照合が不可能)。感覚の不一致が心の状態を作る、という話は \textbf{心理学入門}(知覚と感情)で体系的に学べます。

\section{乗り物のなぜ}

\subsection{飛行機はなぜ飛ぶ}

\textbf{短い答え}: 翼が空気を下に押しやり、その反作用で翼が上に押されるから。

\textbf{しくみ}: 翼には迎え角がついていて、進むと空気を斜め下に押し下げます。第 3 法則(作用・反作用)により、空気は翼を上向きに押し返す。これが揚力です。翼の上面の流速が速くなり圧力が下がる(ベルヌーイの効果)ことも揚力に寄与しますが、\textbf{「翼が空気を押し下げている」}という絵が一番素直な理解です。揚力は速度の 2 乗に比例するので、離陸には十分な滑走速度が要ります。

\textbf{応用と深掘り}: 高速で飛ぶ飛行機の翼端に白い渦(翼端渦)が見えるのは、押し下げた空気の名残です。鳥の V 字編隊は、この渦を隣の鳥が有効利用して省エネしている、という知見があります(生き物のなぜの節)。力の釣り合いと第三法則は \textbf{高校物理やり直し}(力と運動)で、流体の圧力は \textbf{機械のしくみ入門}(ポンプの節)で扱います。

\subsection{エスカレーターはなぜ遅い}

\textbf{短い答え}: 危険の増加が速さの利益を上回るから(速度の 2 乗で事故が重くなる)。

\textbf{しくみ}: エスカレーターの事故は、転倒・巻き込み・落下です。\textbf{乗客の歩く速度 + エスカレーターの速度}の合成で人が動くので、設備速度を 2 倍にすると、歩いた人間の速度も 2 倍になり、\textbf{転倒時の衝撃は 4 倍}(運動エネルギーは速度の 2 乗)になります。高齢者や子どもの歩行速度との差も大きくなる。輸送力の利益より安全コストが大きい、と判断されているのが実情です。

\textbf{応用と深掘り}: 「速度の 2 乗でエネルギーが増える」は、自動車の安全設計(衝突速度と被害)にもそのまま通用します。運動エネルギー \( \frac{1}{2}mv^{2} \) の話は \textbf{高校物理やり直し}(仕事とエネルギー)の中心です。

\subsection{磁気浮上はなぜ浮く}

\textbf{短い答え}: 磁石の反発(または誘導される電流の反発)で、重力と釣り合う力を作るから。

\textbf{しくみ}: リニアモーターカーの浮上には 2 方式があります。\textbf{反発方式}は、車体の超電導磁石と地上コイルの同じ極の反発で浮きます(浮上の高さは安定しない)。\textbf{吸引方式}は、車体をレールに引き寄せて浮かせ、隙間センサーで距離を保ちます。いずれも\textbf{車輪がないので摩擦がなく、速度と騒音の制約が激減する}のが利点です。

\textbf{応用と深掘り}: 磁石の力の正体は電流と磁場(第 4 節のモーターと同じ)で、リニアの推進もモーターを地上に展開したものです。磁場と力の話は \textbf{高校物理やり直し}(電気と回路)と \textbf{機械のしくみ入門}(モーターの節)で、超電導の話は \textbf{量子力学の直感}(物性の入口)につながります。

\section{宇宙と地球のなぜ}

\subsection{月はなぜ同じ面しか見せない}

\textbf{短い答え}: 月の自転周期と公転周期がちょうど等しい(潮汐固定)から。

\textbf{しくみ}: 月は地球の周りを約 27.3 日で公転し、\textbf{同じ 27.3 日で自転}しています。だから地球からは常に同じ面が見えます。偶然ではなく、地球の潮汐力が月の変形と回転を長い時間をかけて調整し、\textbf{安定な状態に落ち着いた}結果です。裏側は 1959 年にソ連の探査機が初めて撮影するまで、誰も見たことがありませんでした。

\textbf{応用と深掘り}: 同じ現象は、人工衛星にも使われています。通信衛星を「常に同じ地点の上空に居させる」(静止軌道)には、公転周期を地球の自転と合わせます。\textbf{周期の同期}という考え方は、潮汐固定と静止軌道で同じです。万有引力と円運動は \textbf{高校物理やり直し}(力と運動)で、周期の数学は \textbf{高校数学やり直し II}(三角関数)で学べます。

\subsection{季節はなぜ生まれる}

\textbf{短い答え}: 地軸が 23.4 度傾いていて、季節によって太陽光の当たり方が変わるから。

\textbf{しくみ}: よくある誤解は「夏は太陽に近いから」です。地軸が傾いているため、夏は\textbf{太陽が高く昇り、日照時間が長い}。同じ面積の地面に受ける光のエネルギーが増え(角度が浅いと同じ光が広い面積に散る)、昼が長いぶん合計の受熱も増えます。南半球では季節が逆になるのも、傾きの証拠です(地球と太陽の距離は 1 年でわずかしか変わりません)。

\textbf{応用と深掘り}: 「角度が浅いと同じ光が広がる」は、灯りの傾け方で部屋の明るさが変わるのと同じ幾何です。太陽光発電のパネルが角度を追うのも同じ理由です。地球と太陽の話は \textbf{機械のしくみ入門}には収めていませんが、\textbf{高校物理やり直し}(円運動と重力)の延長線上にあります。

\subsection{星はなぜ瞬く}

\textbf{短い答え}: 星の光が大気の揺らぎでわずかに曲げられ、明るさがゆらぐから。

\textbf{しくみ}: 大気は温度の違う空気の層が混ざり合っており、それぞれ屈折率が微妙に違います。星の光はこの揺らぐ層を通るたびにわずかに曲げられ、地上では明るさと位置がゆらぎます(シンチレーション)。地球の大気の外(宇宙空間)では瞬かない。惑星が瞬かないのは、\textbf{星より見かけの大きさがずっと広い}ので、揺らぎが全体では平均化されるからです。

\textbf{応用と深掘り}: 宇宙望遠鏡を大気の上に置く理由は、この揺らぎを避けるためです。地上の望遠鏡でも、大気の揺らぎをリアルタイムで補正する技術(補償光学)があります。観測と誤差の話は \textbf{統計学入門}(測定と推定)、屈折の基礎は \textbf{高校物理やり直し}(音と光)です。

\section{生き物のなぜ}

\subsection{葉はなぜ緑}

\textbf{短い答え}: 葉が赤と青の光を吸収して使うので、残りの緑が反射して見えるから。

\textbf{しくみ}: 植物の葉緑素(クロロフィル)は、光合成でエネルギーを集めますが、\textbf{太陽光の全色を同じように使えません}。赤と青の波長を特に良く吸収し、緑をあまり吸収しない。使わない緑が反射して目に届くので、葉は緑に見えます。なぜ緑を余らせるのか(海中の祖先の光環境などの仮説がありますが)は、進化の未解決問題です。

\textbf{応用と深掘り}: 植物工場の栽培照明が「赤と青の LED」で作られているのは、葉がよく吸収する波長に絞った省エネ設計です。光の波長と吸収の話は \textbf{高校化学やり直し}(粒子モデル)、光の性質は \textbf{高校物理やり直し}(音と光)で学べます。

\subsection{鳥はなぜ V 字で飛ぶ}

\textbf{短い答え}: 先頭の翼端にできる上向きの渦を、後ろの鳥が利用して省エネできるから。

\textbf{しくみ}: 翼は空気を押し下げている(乗り物のなぜの節)ので、翼端からは外へ逃げる上向きの渦が出ます。後続の鳥が\textbf{前の鳥の翼端渦の上昇流の斜め外側}に位置すると、自分の翼が揚力を追加でもらえます。研究では編隊飛行でかなりの省エネになると推定されており、群れは疲労を分け合うためにも交代します。

\textbf{応用と深掘り}: 同じ渦の話は、旅客機の離陸間隔が空港で制限される理由(先行機の翼端渦は後続機に危険)でもあります。流れと力の話は \textbf{高校物理やり直し}(力と運動)、群れの協力の話は \textbf{ゲーム理論入門}(繰り返しゲーム)と接続します。

\subsection{猫はなぜ箱に入る}

\textbf{短い答え}: 狭い囲いが、捕食者から隠れられて攻撃を仕掛けられる安心の場所だから。

\textbf{しくみ}: 猫は中間の捕食者です(獲物にもなり、狩りもする)。狭窄した場所は\textbf{背中を守り、出入口を監視できる}戦略的な位置で、休息中の警戒コストが下がります。箱だけでなく、輪の上、段ボールの四角い跡の上(実験でも確認)に座るのは、\textbf{境界が視覚的に示されただけでも安心効果がある}という知見です。

\textbf{応用と深掘り}: 「安心は囲いから生まれる」は人間にも通じる話で、人間の行動の「場の力」の議論と対比できます。\textbf{心理学入門}(発達の安全基地・社会心理)は、この話の科学版です。動物の行動は進化の視点(\textbf{高校生物の分野})でも読めます。

\section{数字と単位のなぜ}

\subsection{1 ダースはなぜ 12}

\textbf{短い答え}: 12 は約数が多い(1, 2, 3, 4, 6)ので、分けるときに扱いやすいから。

\textbf{しくみ}: 物々交換の時代、卵やパンは\textbf{半分に分ける、3 つに分ける、4 つに分ける}場面が頻出しました。10 は約数が 1, 2, 5 しかなく、3 等分できません。12 は 2・3・4・6 で割り切れるので、多くの文化で「まとめ売りの単位」になりました。ダース(12)・グロス(144 = 12 × 12)の体系は、この利便性の名残です。

\textbf{応用と深掘り}: 約数の多さを基準に数を選ぶ発想は、現在でも時間(12 時間・60 分)や角度(360 度 = 約数が大量)に残っています。約数と場合の数は \textbf{高校数学やり直し}(集合と場合の数)で、\textbf{暗号の数学入門}(素因数分解)では素数と約数の話が鍵になります。

\subsection{時間はなぜ 60 進法}

\textbf{短い答え}: 古代バビロニアが 60 を基数に使っていた(約数が多く、1 から数えやすい)から。

\textbf{しくみ}: 60 = 12 × 5 で、2, 3, 4, 5, 6, 10, 12, 15, 20, 30 で割れます。\textbf{指で数える文化}(片手 5 本 × 12 = 60、親指で反対の手の節を数える方法)と結びつき、天文学の角度と時間の単位として 4000 年残りました。メートル法や 10 進法が普及しても、時間と角度だけが 60 進法のまま残ったのは、\textbf{世界中の暦と時計が既に 60 で統一されていた}からです。

\textbf{応用と深掘り}: 60 進法は「割りやすい数」を求めた結果の文化遺産です。同じ発想の現代版が、コンピュータの 2 進法(2 で割り続けられる・電気の ON/OFF に合う)です。進法と素因数分解の話は \textbf{暗号の数学入門}(モジュラ計算)で、データの表現は \textbf{Web フロントエンド入門}(数値の節)にあります。

\subsection{郵便番号はなぜ数字の塊}

\textbf{短い答え}: 地域を階層的に分類するための、自動仕分けのためのコードだから。

\textbf{しくみ}: 日本の郵便番号(例 123-4567)は、上位 3 桁が大きな地域、下 4 桁が町や街区まで細分します。\textbf{番号の先頭から意味が読める}階層構造にすることで、機械が数字を読むだけで配達区分まで決められます。人間の住所の書き方(都道府県 → 市 → 町)と同じ順序を、固定長の数字に翻訳したものです。

\textbf{応用と深掘り}: 階層コードの発想は、電話番号・ ISBN・コンピュータの IP アドレスまで広く使われています。「分類を番号に翻訳する」は、\textbf{データベース入門}(SQL)と \textbf{アルゴリズムとデータ構造}(検索の速さ)で学べます。

\section{まとめ — 雑学から本編へ}

8 つのテーマを振り返ると、雑学の正体は\textbf{既に学んだ道具の応用例}でした。

\begin{itemize}
\item 空の色・虹・雪の結晶 → 波長、屈折、結晶(物理と化学の道具)
\item 塩と雪・吹きこぼれ・冷凍焼け → 溶液、沸点、状態変化(化学の道具)
\item 雷・サイレン・無響室 → 波の速さ、ドップラー、反射(物理の道具)
\item 指紋・聴力・めまい → センサーと感覚(機械と心理学の道具)
\item 飛行機・エスカレーター・リニア → 作用反作用、エネルギー、磁場(物理の道具)
\item 月・季節・星の瞬き → 公転、地軸、大気(物理と地球の道具)
\item 葉の緑・V 字・猫の箱 → 光の吸収、空力、行動(化学・物理・心理学)
\item 12 と 60・郵便番号 → 約数、進法、階層分類(数学と情報の道具)
\end{itemize}

雑学の学び方で一番効くのは、\textbf{気になった「なぜ」を自分の言葉で 1 行で答えてみること}です。1 行で詰まったら、この本で該当の節を読み返し、それでも足りなければ「続きはこちら」の本編に進みます。雑学は入口であって、理解は本編で完成します。

この本は図書館のすべての棚に入口を置きました。読み返すときは、\textbf{興味のある節だけ}を拾い読みしてください。各節の「続きはこちら」が、あなたの次の教科書です。

\end{document}
